% --- ARCHIVO: Capitulos/elementos.tex ---

\chapter{Los Cuatro Elementos de Keitros}

\begin{loreblock}
``Antes de que existiera la primera espada, antes de que el primer puño se cerrara con rabia — existían los cuatro. No como fuerzas ciegas, sino como voluntades que aún no habían encontrado quién las portara.''\\[4pt]
\hfill — \textit{El Moderador, fragmento del Archivo No.~0}
\end{loreblock}

\vspace{0.4cm}

\section{La Naturaleza de los Elementos}

El mundo de Keitros no se rige por los elementos simples que otras civilizaciones conocen: no hay fuego, ni agua, ni tierra, ni aire como fuentes de poder. En Keitros, la energía que mueve el mundo, que da vida, que mata, que destruye y crea — tiene cuatro formas fundamentales.

Estas cuatro formas son el \textbf{Éter}, el \textbf{Vacío}, la \textbf{Luz} y la \textbf{Oscuridad}.

No son fuerzas en equilibrio perfecto. Son fuerzas que \textit{intentan} estarlo. Y en ese intento constante, se libra cada batalla en este mundo.

\begin{imagebox}{Diagrama cósmico de los cuatro elementos — sus relaciones, sus colores simbólicos y las facciones que los representan: Éter (Evol), Vacío (Biskrorum), Luz (Táin), Oscuridad (Déin)}
    \vspace{4cm}
\end{imagebox}

% ──────────────────────────────────────────────────────────────────────────────
\section{El Éter — El Código del Mundo}

\begin{tcolorbox}[enhanced, colback=DeepBlack!85, colframe=GoldRule,
    title={\color{GoldRule}\textbf{EL ÉTER}},
    fonttitle=\large\bfseries,
    attach boxed title to top left={xshift=4mm, yshift=-2mm},
    boxed title style={colback=DeepBlack, colframe=GoldRule}]
    \color{white}
    \itshape
    ``El Éter no se ve. No se toca. No se escucha. Y sin embargo, está en todas partes, en la vibración de cada átomo, en el espacio entre la materia y el vacío. Solo aquellos que aprenden a leerlo pueden escribir en él.''\\[4pt]
    \hfill — \textit{Archimago Vaërzär, La Gran Corte Suprema de la Evolgenia}
\end{tcolorbox}

\vspace{0.3cm}

El \textbf{Éter} es el elemento más complejo y el más estudiado. Está presente en absolutamente todo — en el aire, en la piedra, en la sangre de los seres vivos, en el espacio entre las estrellas. Sin embargo, la mera presencia del Éter no implica su dominio: el Éter en estado ``sin uso'' es energía potencial pura, invisible e inerte.

Fueron los \textbf{hijos de Evol} quienes descubrieron que el Éter puede ser transformado. A través del correcto entendimiento de sus frecuencias y patrones — lo que la Evolgenia llama \textit{``lectura del código''} — es posible convertir el Éter latente en Éter activo: una fuerza capaz de modificar la materia, alterar estructuras físicas y manifestar cambios en la realidad observable.

Esta es la base de toda la ciencia y magia de los Evolgens: no crean energía, \textbf{la reescriben}.

\begin{reglaespecial}{Éter en Juego}
\begin{itemize}
    \item El Éter \textbf{no interfiere} con la Luz ni con la Oscuridad. Sus efectos son independientes de estos elementos.
    \item Las Zonas Secas representan zonas donde el Éter ha sido agotado temporalmente por uso excesivo.
    \item Los Evolgens pueden detectar concentraciones de Éter en el entorno con tiradas de Focus.
    \item Habilidades de tipo \textbf{``Modificación''} y \textbf{``Transmutación''} usan Éter como fuente de energía.
\end{itemize}
\end{reglaespecial}

% ──────────────────────────────────────────────────────────────────────────────
\section{El Vacío — La Ausencia que Pesa}

\begin{tcolorbox}[enhanced, colback=DeepBlack!95, colframe=VoidPurple!80,
    title={\color{VoidPurple!80}\textbf{EL VACÍO}},
    fonttitle=\large\bfseries,
    attach boxed title to top left={xshift=4mm, yshift=-2mm},
    boxed title style={colback=DeepBlack, colframe=VoidPurple}]
    \color{white}
    \itshape
    ``Cuando los demás ven nada, nosotros vemos todo. El vacío no es ausencia — es el estado puro de la posibilidad antes de que la materia decida existir.''\\[4pt]
    \hfill — \textit{Thothk'thun, Rey de Abyssis}
\end{tcolorbox}

\vspace{0.3cm}

El \textbf{Vacío} es un elemento que, al igual que el Éter, está presente en el entorno — pero a diferencia de este, es casi inaccesible para la gran mayoría de las especies de Keitros. Solo los \textbf{hijos de Biskrorum}, criaturas nacidas de un ser interdimensional cuyo origen es el vacío estelar, pueden percibir, canalizar y utilizar esta energía de forma natural.

El Vacío no es simplemente la ausencia de materia. Dentro de un entorno con materia, el Vacío actúa como una fuerza que \textbf{reescribe las leyes físicas locales}: puede anular la gravedad en un radio determinado, crear zonas donde la luz no entra ni sale, doblar el espacio para acelerar o detener el movimiento. Para los Biskrorum, el Vacío es su lengua materna — ven el mundo a través de sus reglas, y sus sentidos están naturalmente sintonizados con él.

Para cualquier otra especie, el Vacío es imperceptible y en muchos casos \textbf{aterrador}: un campo de Vacío activo puede producir desorientación, pérdida del sentido del espacio y pánico instintivo en criaturas no adaptadas.

\begin{reglaespecial}{Vacío en Juego}
\begin{itemize}
    \item El Vacío \textbf{no interfiere} con la Luz ni con la Oscuridad.
    \item Las unidades no-Biskrorum que terminen su turno dentro de un campo de Vacío activo sufren \textbf{--1 a Focus} y deben tirar Res. Mágica para no quedar \textbf{Desorientadas} (pierden la orientación de dirección — el oponente elige hacia dónde miran al inicio de su turno).
    \item Las habilidades de Vacío pueden crear zonas de \textbf{gravedad alterada}, \textbf{silencio absoluto} (impiden Hechizos Rápidos) o \textbf{distorsión espacial} (modifica distancias de MOV).
    \item El Vacío \textbf{sí choca con el Éter}: cuando ambas energías se encuentran en el mismo espacio, se anulan mutuamente en una explosión de energía pura (daño a todas las unidades en el área).
\end{itemize}
\end{reglaespecial}

% ──────────────────────────────────────────────────────────────────────────────
\section{La Luz — El Alma del Mundo}

\begin{tcolorbox}[enhanced, colback=DeepBlack!80, colframe=GoldRule!60,
    title={\color{GoldRule!90}\textbf{LA LUZ}},
    fonttitle=\large\bfseries,
    attach boxed title to top left={xshift=4mm, yshift=-2mm},
    boxed title style={colback=DeepBlack, colframe=GoldRule!60}]
    \color{white}
    \itshape
    ``No hablamos del sol. Hablamos de aquello que hace que la mirada de un ser vivo tenga peso. La Luz es la firma individual de cada alma — el patrón único que el Gran Árbol reconoce como 'uno de los suyos'.''\\[4pt]
    \hfill — \textit{Ebánory, el Héroe Semilla, registro del Gran Árbol}
\end{tcolorbox}

\vspace{0.3cm}

La \textbf{Luz} no es la luz solar. No es luminosidad, ni calor, ni fotones. La Luz es la \textbf{energía vital de las almas} — la firma energética única que hace que un ser esté vivo, consciente, presente. Fue Táin, el gran árbol del mundo, quien trajo este elemento a Keitros: sus semillas, al dar vida a los Fäes y a los Elementales, introdujeron almas en el mundo, y con ellas, la Luz.

Los \textbf{hijos de Táin} pueden ver, tocar y manipular esta energía como herramienta. Pueden drenarla de enemigos para debilitarlos o matarlos, pueden concentrarla para sanar a aliados, pueden canalizarla para potenciar el crecimiento de plantas y terreno, o pueden proyectarla como un arma de energía pura que quema a los no-vivos.

Las almas individuales brillan con Luz. Las grandes concentraciones de seres vivos en un territorio generan más Luz en el entorno. Y cuando un ser muere bajo la influencia de un hijo de Táin, su Luz puede ser absorbida y redistribuida.

\begin{reglaespecial}{Luz en Juego}
\begin{itemize}
    \item La Luz \textbf{no es afectada} por entornos ricos en Éter o Vacío.
    \item La Luz \textbf{choca directamente con la Oscuridad}: cuando ambas se encuentran, se produce una reacción de aniquilación. Ver sección de Relaciones entre Elementos.
    \item Los hijos de Táin pueden detectar unidades vivas \textbf{a través de la cobertura} si tienen una habilidad de percepción de Luz (no requiere Line of Sight para detectar — pero sí para atacar).
    \item Las habilidades de Luz pueden \textbf{sanar}, \textbf{drenar vitalidad}, \textbf{crear barreras de alma} o \textbf{revelar unidades ocultas} dentro del radio.
    \item Los \textbf{No-Vivos} (esqueletos, autómatas, construcciones) son inmunes a los efectos de drenaje de Luz pero reciben daño doble de ataques de Luz que sean ``Ofensivos''.
\end{itemize}
\end{reglaespecial}

% ──────────────────────────────────────────────────────────────────────────────
\section{La Oscuridad — El Peso de los Vivos}

\begin{tcolorbox}[enhanced, colback=DeepBlack!95, colframe=HammerRed!70,
    title={\color{HammerRed!80}\textbf{LA OSCURIDAD}},
    fonttitle=\large\bfseries,
    attach boxed title to top left={xshift=4mm, yshift=-2mm},
    boxed title style={colback=DeepBlack, colframe=HammerRed}]
    \color{white}
    \itshape
    ``¿Crees que la oscuridad es ausencia de luz? La oscuridad es lo que dejás cuando sufrís. Es lo que exhalás cuando el miedo te aprieta el pecho. Nosotros no nos alimentamos de sombras — nos alimentamos de vos.''\\[4pt]
    \hfill — \textit{César, el Gran Héroe Daemon, al ejército capturado del Instituto del Progreso}
\end{tcolorbox}

\vspace{0.3cm}

La \textbf{Oscuridad}, al igual que la Luz, no es lo que parece. No es la ausencia de iluminación, no es la noche ni las sombras. La Oscuridad es una energía que se \textbf{genera en los seres vivos}: cada vez que un ser siente miedo, dolor, odio, desesperación o cualquier emoción negativa intensa, libera Oscuridad al mundo. Esta energía se acumula en el entorno de forma natural y permanente.

Los espacios de sufrimiento histórico — campos de batalla, mazmorras, lugares de tortura, ciudades arrasadas — están cargados de Oscuridad incluso siglos después. Los \textbf{hijos de Déin} se alimentan de esta energía para vivir, y la canalizan para potenciar sus habilidades con una eficacia directamente proporcional al sufrimiento del entorno en que combaten.

Más aún: los Daemons pueden \textit{inducir} sufrimiento en sus enemigos para generar Oscuridad fresca. El miedo es literalmente combustible.

\begin{reglaespecial}{Oscuridad en Juego}
\begin{itemize}
    \item La Oscuridad \textbf{no es afectada} por entornos ricos en Éter o Vacío.
    \item La Oscuridad \textbf{choca directamente con la Luz}: cuando ambas se encuentran, se produce una reacción de aniquilación. Ver sección siguiente.
    \item Los hijos de Déin recuperan \textbf{PE adicionales} al inicio de cada ronda si hay unidades Enemigas con estados negativos activos (miedo, dolor, agonía).
    \item En territorios con el marcador \textbf{``Zona Oscura''} (creado por ciertas habilidades), los Daemons ganan +1 a todos sus atributos mientras permanezcan en la zona.
    \item Las habilidades de Oscuridad pueden \textbf{infundir miedo} (estado: Aterrado), \textbf{drenar voluntad}, \textbf{reanimar caídos} o \textbf{crear zonas de sombra} que bloquean la Línea de Visión para enemigos no-Daemons.
\end{itemize}
\end{reglaespecial}

% ──────────────────────────────────────────────────────────────────────────────
\section{Relaciones entre Elementos}

Los cuatro elementos no interactúan todos entre sí. Las relaciones son específicas y simétricas:

\begin{imagebox}{Esquema de las relaciones entre elementos — dos pares: [Éter ↔ Vacío] y [Luz ↔ Oscuridad], con las dos líneas de choque marcadas en rojo}
    \vspace{3cm}
\end{imagebox}

\subsection*{Éter y Vacío — La Colisión de Paradigmas}

El Éter y el Vacío son fuerzas opuestas en naturaleza: el Éter transforma la materia, el Vacío reescribe las leyes que la gobiernan. Cuando ambas energías se concentran en el mismo punto, la tensión es insostenible y produce una \textbf{explosión de liberación energética}.

\begin{tcolorbox}[colback=HammerRed!8, colframe=HammerRed, title=\textbf{REGLA: Colisión Éter–Vacío}]
\small Si una habilidad de Éter y una habilidad de Vacío afectan simultáneamente el mismo punto del tablero (o se solapan en su área de efecto), ambas se anulan mutuamente y todas las unidades en el área de solapamiento reciben \textbf{1d6 daño directo} (ignora RD). Coloca un marcador de \textbf{Zona Inestable} en ese punto — mientras dure la partida, cualquier uso de PE en esa zona tiene un 50\% de probabilidad de fallar.
\end{tcolorbox}

\subsection*{Luz y Oscuridad — La Guerra Eterna}

La Luz y la Oscuridad son las dos fuerzas que más directamente se oponen en Keitros. No se anulan entre sí de forma explosiva — se \textbf{consumen}. Cuando ambas se encuentran, la más intensa devora a la más débil.

\begin{tcolorbox}[colback=HammerRed!8, colframe=HammerRed, title=\textbf{REGLA: Consumo Luz–Oscuridad}]
\small Cuando una habilidad de Luz y una de Oscuridad afectan la misma unidad en el mismo turno, comparan sus potencias (el valor de PE invertido en cada una). La de mayor potencia prevalece; la de menor potencia se cancela completamente. En caso de empate exacto, ambas se anulan.\\[4pt]
Esto se aplica también a efectos de zona: un \textbf{Aura de Luz} activa cancela la Oscuridad generada por estados negativos dentro de su radio (los Daemons dentro del aura no pueden recuperar PE por sufrimiento enemigo).
\end{tcolorbox}

\subsection*{Elementos que No Interactúan}

El Éter no afecta ni es afectado por la Luz ni por la Oscuridad — y viceversa. El Vacío tampoco interfiere con la Luz ni con la Oscuridad. Esto significa que en una batalla entre Biskrorum (Vacío) y Táin (Luz), ninguno puede usar su elemento para contrarrestar directamente al otro: deben vencerse con acero, estrategia y unidades.

\begin{center}
\begin{tabularx}{0.85\linewidth}{|c|c|c|c|c|}
\hline
\rowcolor{HammerRed!20} & \textbf{Éter} & \textbf{Vacío} & \textbf{Luz} & \textbf{Oscuridad} \\ \hline
\textbf{Éter}      & —           & 💥 Colisión & ✓ Sin efecto & ✓ Sin efecto \\ \hline
\textbf{Vacío}     & 💥 Colisión & —           & ✓ Sin efecto & ✓ Sin efecto \\ \hline
\textbf{Luz}       & ✓ Sin efecto & ✓ Sin efecto & —           & 🔥 Consumo \\ \hline
\textbf{Oscuridad} & ✓ Sin efecto & ✓ Sin efecto & 🔥 Consumo  & — \\ \hline
\end{tabularx}
\end{center}
